\documentclass[instructions.tex]{subfiles}
\begin{document}
 
\chapter{Des divertissements du monde.}
 
S'il faut quelque récréation et quelque divertissement à l'homme, il doit les prendre selon Dieu; mais les grands repas, les jeux immodérés, les bals et les spectacles, qui sont les divertissemens ordinaires du monde, sont-ils selon les desseins de Dieu? Quatre chefs qu'il est important d'examiner.
 
\primo Les repas. On peut donner quelques repas, quelques festins, selon sa condition : la société et la bienséance l'exigent en certaines conjonctures. Le Sauveur assista lui-même aux festins qu'on lui prépara chez Simon le pharisien, chez Zachée, chez Matthieu; mais les festins doivent être rares et modérés.
 
Les grands repas fréquens et somptueux sont l'écueil de l'honnêteté et de la modération; souvent on y dit ce qu'on n'aurait osé penser à jeun. Il est difficile de n'y point passer les bornes que la tempérance et la religion prescrivent : l'âme s'y ressent de la dissolution du corps; le coeur s'y dissipe; l'esprit de Dieu s'y éteint; on y scandalise les pauvres qui meurent de faim, tandis que d'autres sont trop rassasiés; on y oublie la pudeur et le salut, on y altère sa santé; on y avance sa mort, selon le proverbe : La gourmandise et la table en ont plus fait mourir que le glaive. Au contraire, la sobriété et l'abstinence prolongent les jours\footnote{In multis enim escis erit infirmitas.... Qui autem abstinens est, adjiciet vitam. Eccl. 37.}.
 
\secundo Les jeux. On permet des jeux innocens et des amusemens aux enfans : cela convient à leur âge, mais qu'il sied mal à des gens dont la vie se passe dans l'inutilité, de dire qu'ils jouent par délassement et pour s'amuser! comme si des fainéans et des fainéantes, qui coulent leurs jours dans la mollesse, dont la vie est une perpétuelle récréation, avoient besoin de délassement, et comme si le temps qu'on doit ménager pour le salut et pour ses affaires, n'étoit donné de Dieu que pour se divertir et s'amuser.
 
Autrefois le sénat de Rome, pour honorer la mémoire d'une dame illustre, lui érigea une statue qui tenoit une quenouille, pour montrer qu'elle s'étoit occupée, non pas aux jeux et aux spectacles, mais au soin de sa famille et à des ouvrages convenables à son état et à son sexe, et pour apprendre à toutes les dames que l'attention à leur famille, l'aiguille et le fuseau, sont leurs plus nobles occupations. Ce sont celles que le Saint-Esprit donne à la femme forte\footnote{Digit ejus apprehenderunt fusum. Prov. 31.}. Mais aujourd'hui, si l'on érigeoit des statues à certaines dames, il faudroit leur mettre en main l'attirail d'une toilette, des romans et des cartes.
 
Le délassement et quelques récréations conviennent à des personnes d'étude, à des gens de travail, à des convalescens, comme jeux d'esprit, jeux d'exercice, et autres jeux honnêtes. Les jeux de main ne conviennent point à des gens raisonnables. Jouez peu de temps, peu ou point d'argent. Un jeu trop intéressé, ou qui sert d'occupation, n'est plus innocent.
 
Voilà les règles que suivent les personnes sages, mais il est rare qu'on s'en tienne là. On joue long-temps; on joue sans modération; on joue par intérêt; et de là combien de désordres! ce qui a fait dire à un célèbre magistrat, que les grands joueurs réunissent en eux les grands vices. En voici le détail :
 
Imprudens, ils risquent étourdiment la perte de leur bien sur un coup de dés, sur une carte. Orgueilleux, ils se piquent de savoir jouer; fiers jusqu'à l'impudence quand ils gagnent. L'avarice, l'esprit d'intérêt qui les possède, leur fait employer la tromperie et la surprise. L'envie les tourmente; ils ne voient qu'avec chagrin le bonheur des autres, cachant, sous une contenance forcée, le dépit qui les ronge.
 
D'autres fois opiniâtres, ils ne veulent ni céder, ni entendre raison. Tantôt emportés jusqu'à la fureur et au désespoir, ils blasphêment et s'en prennent à Dieu quand ils perdent. Tantôt injustes et voleurs, jouant au préjudice de leurs créanciers, ou ce qu'ils ont dérobés dans leur famille, ou sur les emprunts qu'ils ne peuvent rendre; enfin, soupçonneux, défis, brouillons, chicaneurs, querelleurs, voilà les passions qui tour-à-tour ou tout à la fois les agitent.
 
Jérémie, ce saint prophète, et la chaste Sara, épouse de Tobie, pour obtenir miséricorde, disaient à Dieu : Seigneur, ayez égard que je n'ai jamais été dans les assemblées des joueurs\footnote{Jer. 15. Tob. 3.}. Vraiment un joueur est bien en état de parler à Dieu, en sortant du jeu avec une conscience agitée de tant de passions! Comment gémiroit-il devant Dieu sur l'état de son âme, lui qui ne daigne pas même y réfléchir, et qui se croit innocent!
 
Ce n'est pas tout. Quelle vertu voit-on dans un joueur habituel? Sans patience, il prend feu pour rien: sans humilité, il veut toujours avoir droit et jamais tort; sans modération, il se choque d'une parole échappée; sans piété, point de retour à Dieu, point d'assistance ni de recueillement au service divin. Dieu veuille que la pudeur soit épargnée dans le jeu, et que la chasteté n'en soit pas la victime! Enfin, sans attention à tous ses devoirs, un joueur néglige sa famille, ses affaires, les sacremens et son salut.
 
Si je n'ai pas la charité, dit saint Paul, je ne suis rien. Un homme passionné pour le jeu n'a pas même l'ombre de cette vertu. Il est sans charité pour sa femme et ses enfans, qu'il réduit à la misère par son acharnement au jeu; sans charité pour son adversaire, qu'il voudroit dépouiller et que souvent il ruine; sans charité pour les jeunes gens, qu'il séduit, qu'il attire au jeu, et qu'il met dans la disgrâce de leurs parens; sans charité pour le pauvre, à qui il refuse un sou, tandis que, etc.; enfin, sans charité pour soi-même, se souciant peu d'altérer sa santé et de perdre son âme, pourvu qu'il se satisfasse et qu'il gagne de l'argent.
 
Oh! combien doit-on craindre cette passion! elle est d'autant plus à craindre, qu'outre les désordres auxquels elle entraîne, on ne se corrige, pour l'ordinaire, de la passion du jeu, que quand on est ruiné, après en avoir ruiné d'autres, et très-souvent la conscience embarrassée de plusieurs restitutions à faire, etc.
 
\section{Résolutions}
 
\primo Je ne me laisserai pas entraîner à l'amour du jeu.

\secundo Je ne me permettrai aucun jeu de hasard; je jouerai peu et je m'abstiendrai de tout jeu de mains qui pourrait devenir dangereux à la vertu.
 
\prayer{Mon Dieu, faites-moi la grâce d'user toujours de modération dans les récréations que je suis obligé de prendre, et de m'y comporter avec tant de sagesse, que j'en sorte constamment sans vous offenser.}
 
\end{document}
